\documentclass[a4paper,11pt]{ltxdoc}
\usepackage[UTF8,scheme=plain]{ctex}
\usepackage{geometry}
\usepackage{xcolor}
\usepackage{booktabs}
\usepackage{longtable}
\usepackage{array}
\usepackage{enumitem}
\usepackage{listings}
\usepackage{hypdoc}

\geometry{margin=2.55cm}
\setlength{\emergencystretch}{3em}
\hypersetup{
  colorlinks=true,
  linkcolor=blue!55!black,
  urlcolor=blue!55!black,
  citecolor=blue!55!black,
  pdftitle={Maki Notes Implementation Manual},
  pdfauthor={Maki Notes}
}

\definecolor{MakiBlue}{RGB}{42,91,142}
\definecolor{MakiGray}{RGB}{246,247,249}
\definecolor{MakiRule}{RGB}{210,216,224}

\lstdefinestyle{maki-doc}{
  basicstyle=\ttfamily\small,
  columns=fullflexible,
  keepspaces=true,
  frame=single,
  framerule=0.3pt,
  rulecolor=\color{MakiRule},
  backgroundcolor=\color{MakiGray},
  breaklines=true,
  tabsize=2
}
\lstset{style=maki-doc}

\newcommand{\mfile}[1]{\texttt{#1}}
\newcommand{\mcls}[1]{\textsf{#1}}
\newcommand{\mpkg}[1]{\textsf{#1}}
\newcommand{\mopt}[1]{\texttt{#1}}
\newcommand{\menv}[1]{\texttt{#1}}
\newcommand{\mcmd}[1]{\texttt{\textbackslash#1}}
\newcommand{\mkey}[1]{\texttt{#1}}
\newcommand{\Maki}{\textsf{Maki Notes}}

\newenvironment{implementation}[1]
  {\par\medskip\noindent\textbf{\color{MakiBlue}Implementation note: #1}\par\smallskip}
  {\par\medskip}

\newenvironment{interfaceblock}[1]
  {\par\medskip\noindent\textbf{\color{MakiBlue}Interface: #1}\par\smallskip\begin{description}[leftmargin=3.2cm,style=nextline]}
  {\end{description}\par\medskip}

\title{\Maki{} Implementation Manual\\
  \large A Functional Guide to the \mfile{.cls} and \mfile{.sty} Sources}
\author{Maki Notes Project}
\date{2026-04-18}

\begin{document}
\maketitle
\tableofcontents

\begin{abstract}
This document is an implementation manual for maintainers. It explains how
\mfile{maki-notes.cls}, \mfile{maki-beamer.cls}, \mfile{maki-notes.sty}, and
\mfile{beamerthememaki.sty} work internally. The structure follows the style of
professional CTAN package documentation: file map first, public interfaces next,
then implementation notes grouped by function. It covers option parsing, mode
detection, font and theme selection, metadata workflows, theorem-like boxes,
solution modes, cross-reference helpers, TikZ semantic styles, Beamer templates,
and build automation.
\end{abstract}

\section{Source File Architecture}

\Maki{} is organized as two document-class entry points and two style layers.

\begin{longtable}{@{}p{0.26\linewidth}p{0.68\linewidth}@{}}
\toprule
File & Responsibility \\
\midrule
\mfile{maki-notes.cls} &
Book-style lecture-note entry point. It parses \mopt{fontpreset}, \mopt{theme},
\mopt{answers}, and \mopt{solutions}; loads \mcls{ctexbook}; initializes cover
macros; and loads \mfile{maki-notes.sty}.\\
\mfile{maki-beamer.cls} &
Presentation entry point. It parses \mopt{layout}, \mopt{plain}, \mopt{invert},
and \mopt{accessibility}; loads \mcls{ctexbeamer}; and forwards theme options to
\mfile{beamerthememaki.sty}.\\
\mfile{maki-notes.sty} &
Shared style implementation. It owns font presets, semantic themes, mathematical
typography, course and research metadata, guide pages, tcolorbox environments,
solution collection, reference helpers, wrap figures, margin notes, and
TikZ/PGFPlots semantic styles.\\
\mfile{beamerthememaki.sty} &
Beamer theme implementation. It defines title pages, outline pages, headlines,
footlines, progress bars, block colors, and three layout variants.\\
\bottomrule
\end{longtable}

\begin{implementation}{Thin entry classes}
The \mfile{.cls} files intentionally remain thin. They parse options, choose the
base class, define a few top-level document commands, and then delegate visual
and semantic behavior to the style layers. This keeps notes and slides aligned.
\end{implementation}

\section{Document Classes and Option Flow}

\subsection{\mfile{maki-notes.cls}}

The notes class uses \mpkg{kvoptions} to declare stable class options:

\begin{lstlisting}[language=TeX]
\DeclareStringOption[common]{fontpreset}
\DeclareStringOption[default]{theme}
\DeclareBoolOption[false]{answers}
\DeclareStringOption[inline]{solutions}
\DeclareDefaultOption{\PassOptionsToClass{\CurrentOption}{ctexbook}}
\ProcessKeyvalOptions*
\end{lstlisting}

\begin{interfaceblock}{Notes class options}
\item[\mopt{fontpreset}] One of \mopt{common}, \mopt{auto}, \mopt{windows}, \mopt{macos}, or \mopt{linux}.
\item[\mopt{theme}] One of \mopt{default}, \mopt{ocean}, \mopt{forest}, \mopt{graphite}, \mopt{amber}, \mopt{berry}, \mopt{sandstone}, \mopt{cobalt}, \mopt{sage}, \mopt{ruby}, or \mopt{midnight}.
\item[\mopt{answers}] Boolean switch controlling whether solution bodies are visible.
\item[\mopt{solutions}] Output mode for solutions: \mopt{inline}, \mopt{appendix}, or \mopt{hidden}.
\end{interfaceblock}

After option processing, the class copies option values into macros consumed by
\mfile{maki-notes.sty}:

\begin{lstlisting}[language=TeX]
\edef\MakiFontPreset{\maki@fontpreset}
\edef\MakiTheme{\maki@theme}
\ifmaki@answers
  \MakiShowAnswerstrue
\else
  \MakiShowAnswersfalse
\fi
\edef\MakiSolutionsMode{\maki@solutions}
\end{lstlisting}

\subsection{\mfile{maki-beamer.cls}}

The Beamer class mirrors the same pattern, but uses \mcls{ctexbeamer} as the base
class and adds presentation-specific options. The selected layout and boolean
theme flags are assembled into \mcmd{MakiBeamerThemeOptions}; the class then
loads the theme by expanding that option list into \mcmd{usetheme}.

\begin{lstlisting}[language=TeX]
\expandafter\usetheme\expandafter[\MakiBeamerThemeOptions]{maki}
\end{lstlisting}

\begin{implementation}{Dynamic theme option forwarding}
The option list is built incrementally because \mopt{layout}, \mopt{plain},
\mopt{invert}, and \mopt{accessibility} are parsed at the class level but applied
by the Beamer theme. Expanding the list before \mcmd{usetheme} keeps the user
interface at the \mcmd{documentclass} level.
\end{implementation}

\section{Shared Style Mode Detection}

\mfile{maki-notes.sty} begins by detecting whether it is running under Beamer:

\begin{lstlisting}[language=TeX]
\newif\ifMakiBeamerMode
\MakiBeamerModefalse
\@ifclassloaded{beamer}{\MakiBeamerModetrue}{}
\@ifclassloaded{ctexbeamer}{\MakiBeamerModetrue}{}
\end{lstlisting}

This switch controls package loading and feature availability:

\begin{itemize}
\item Notes mode loads \mpkg{geometry}, \mpkg{fancyhdr}, \mpkg{caption}, \mpkg{wrapfig}, and \mpkg{marginnote}.
\item Beamer mode skips book-specific page geometry, headers, captions, and marginalia.
\item Shared mathematical environments, research blocks, and TikZ styles remain available.
\end{itemize}

\begin{implementation}{Compatibility boundary}
The shared style is not intended to be a standalone generic package. It is a
support layer for \mcls{maki-notes} and \mcls{maki-beamer}. The mode switch
prevents conflicts between Beamer and book-oriented packages.
\end{implementation}

\section{Font Presets}

Font setup has two layers: Latin fonts and CJK fonts.

\subsection{Latin Fonts}

\mcmd{makisetlatinfonts} prefers TeX Gyre Pagella for text, TeX Gyre Heros for
sans serif, and Latin Modern Mono for monospaced text. These fonts are usually
available in TeX Live installations, making the default path portable.

\subsection{CJK Fonts}

CJK presets are implemented by named dispatcher macros:

\begin{longtable}{@{}p{0.38\linewidth}p{0.54\linewidth}@{}}
\toprule
Macro & Strategy \\
\midrule
\mcmd{makiapplyfontpresetcommon} & Uses Fandol fonts as the safest TeX Live fallback.\\
\mcmd{makiapplyfontpresetwindows} & Detects \mopt{SimSun}, then chooses common Windows CJK fonts.\\
\mcmd{makiapplyfontpresetmacos} & Detects \mopt{Songti SC}, then chooses common macOS CJK fonts.\\
\mcmd{makiapplyfontpresetlinux} & Detects Noto or Source Han CJK fonts, otherwise falls back to \mopt{common}.\\
\mcmd{makiapplyfontpresetauto} & Probes macOS, Windows, and Linux candidates in order, then falls back to \mopt{common}.\\
\bottomrule
\end{longtable}

\begin{implementation}{Font probing}
\mcmd{makipickfirstfont} receives a comma-separated list of candidate font names
and stores the first available one in \mcmd{MakiPickedFont}. If no candidate is
available, it stores a known fallback. Explicit but unavailable presets emit a
\mcmd{PackageWarningNoLine} and fall back to \mopt{common}.
\end{implementation}

\section{Semantic Theme System}

Themes are implemented by macros named \mcmd{makiapplytheme<name>}. Each theme
defines the same semantic color slots:

\begin{itemize}
\item \mkey{LogicColor}/\mkey{LogicBg}: theorem, proof, and logical emphasis.
\item \mkey{DefColor}/\mkey{DefBg}: definitions, axioms, and input/output blocks.
\item \mkey{ExColor}/\mkey{ExBg}: examples, exercises, and primary flow nodes.
\item \mkey{NoteColor}: notes, methods, and research ideas.
\item \mkey{GeometryColor}/\mkey{GeometryBg} and \mkey{TopologyColor}/\mkey{TopologyBg}: geometry and topology diagrams.
\item \mkey{NNInputColor}, \mkey{NNHiddenColor}, \mkey{NNOutputColor}, and \mkey{NNTensorColor}: neural-network diagrams.
\item \mkey{MainText}: main body text color.
\end{itemize}

\begin{implementation}{Why semantic colors matter}
Boxes, TikZ styles, and Beamer templates all consume semantic color slots rather
than hard-coded palette values. A new theme only needs to define the slots and
register its name in the documentation and tests.
\end{implementation}

\section{Mathematical Typography}

\mfile{maki-notes.sty} applies lecture-note-oriented mathematical defaults:

\begin{itemize}
\item \mcmd{ge} and \mcmd{le} are mapped to \mcmd{geqslant} and \mcmd{leqslant}.
\item \mcmd{lim}, \mcmd{int}, and \mcmd{sum} use display style by default.
\item \mcmd{dfrac} adds invisible struts to improve vertical spacing.
\item Matrix and cases environments use LaTeX hooks to increase line spread.
\end{itemize}

\begin{implementation}{Scope of the global math changes}
These changes intentionally optimize lecture-note readability. If the project is
ever split into a generic CTAN package, these global redefinitions should become
optional.
\end{implementation}

\section{Course Metadata and Chapter Guides}

\subsection{Key--Value Interface}

Course, chapter, and guide metadata are parsed with \mpkg{pgfkeys}:

\begin{lstlisting}[language=TeX]
\SetCourseInfo{
  course={Advanced Mathematics},
  term={Spring 2026},
  lecture={Lecture 3},
  date={2026-04-13}
}
\end{lstlisting}

\begin{longtable}{@{}p{0.3\linewidth}p{0.62\linewidth}@{}}
\toprule
Key family & Stored state \\
\midrule
\mkey{/maki/course} & Course, term, audience, instructor, lecture, and date.\\
\mkey{/maki/chapter} & Pending chapter metadata, later bound to the next chapter.\\
\mkey{/maki/guide} & Route, formulas, pitfalls, methods, and exam focus fields.\\
\bottomrule
\end{longtable}

\subsection{Binding Chapter State}

In notes mode, \mcmd{pretocmd} hooks into \mcmd{@chapter} and \mcmd{@schapter}:

\begin{lstlisting}[language=TeX]
\pretocmd{\@chapter}{\MakiBindPendingChapterInfo}{}{}
\pretocmd{\@schapter}{\MakiBindPendingChapterInfo}{}{}
\end{lstlisting}

\mcmd{MakiBindPendingChapterInfo} copies pending fields into current chapter
fields, resolves effective lecture/date values, and clears the pending state.
This lets authors write \mcmd{SetChapterInfo} before a chapter without leaking
metadata into later chapters.

\subsection{Guide Page Generation}

\mcmd{MakeChapterGuide} builds a chapter-opening guide page:

\begin{itemize}
\item The chapter title is captured through \mcmd{chaptermark}.
\item Course information is read from \mcmd{SetCourseInfo}.
\item Learning goals, route, formulas, methods, pitfalls, and exam focus are rendered with \menv{tcbraster}.
\item Local navigation is generated with \mpkg{etoc} and \mcmd{localtableofcontents}.
\end{itemize}

\section{Research Notebook Layer}

The research layer is parallel to the course layer. It provides project metadata,
entry metadata, research guide cards, research objects, and simple indexes.

\begin{interfaceblock}{Research interfaces}
\item[\mcmd{SetResearchInfo}] Stores project, area, author, collaborators, status, creation date, update date, and tags.
\item[\mcmd{SetEntryInfo}] Stores the current entry topic, source, source key, update date, tags, and objective.
\item[\mcmd{MakeEntryGuide}] Prints the current research entry guide.
\item[\mcmd{DeclareSymbol}] Registers a symbol and its explanation.
\item[\mcmd{PrintSymbolIndex}] Prints all registered symbols.
\item[\mcmd{PrintOpenProblemIndex}] Prints registered questions and conjectures.
\end{interfaceblock}

\begin{implementation}{Dynamic entry storage}
Symbol and open-problem indexes use counters plus dynamic control sequence names.
For example, the \(n\)-th symbol is stored in
\verb|\maki@symbol@n@symbol| and \verb|\maki@symbol@n@desc|. Printing iterates
from 1 to the current counter value with \mcmd{@whilenum}.
\end{implementation}

\section{Wrap Figures and Marginal Notes}

In notes mode, the package loads \mpkg{wrapfig}, \mpkg{marginnote}, and
\mpkg{ifoddpage}. It exposes a small public layer:

\begin{interfaceblock}{Figure/text and margin-note interfaces}
\item[\menv{leftfiguretext}] Wrap a figure on the left and text on the right.
\item[\menv{rightfiguretext}] Wrap a figure on the right and text on the left.
\item[\menv{figuretextleft}] Alias-style wrapper around \menv{leftfiguretext}.
\item[\menv{figuretextright}] Alias-style wrapper around \menv{rightfiguretext}.
\item[\mcmd{sidenote}] Untitled marginal note.
\item[\mcmd{marginremark}] Marginal note with an optional title.
\item[\mcmd{sidecaption}] Figure caption in the margin.
\end{interfaceblock}

\begin{implementation}{Outer-margin placement}
\mcmd{maki@outernote} uses \mcmd{checkoddpage}. In two-sided documents it chooses
right-side formatting on odd pages and left-side formatting on even pages; in
one-sided documents it always uses the right-side format.
\end{implementation}

\section{TikZ and PGFPlots Semantic Styles}

TikZ styles are grouped by meaning rather than by raw appearance.

\begin{longtable}{@{}p{0.29\linewidth}p{0.63\linewidth}@{}}
\toprule
Style family & Use case \\
\midrule
\mkey{nn ...} & Neural-network neurons, tensors, convolution blocks, skip links, and attention edges.\\
\mkey{transformer ...} & Tokens, attention blocks, FFNs, norms, residual flows, and memory nodes.\\
\mkey{pinn ...}/\mkey{fpinn ...} & PINN/fPINN physics constraints, boundary/initial blocks, losses, and fractional operators.\\
\mkey{math ...}/\mkey{geom ...} & Axes, curves, tangents, normals, geometric points, segments, and angle marks.\\
\mkey{topo ...}/\mkey{manifold ...}/\mkey{sphere ...} & Fundamental domains, manifolds, local charts, tangent planes, and spheres.\\
\mkey{fourier ...} & Time-domain signals, frequency-domain spectra, windows, transform arrows, and markers.\\
\mkey{flow ...}/\mkey{timeline ...} & Flowcharts and timelines.\\
\mkey{prob ...}/\mkey{graph ...} & Probability diagrams and graph structures.\\
\bottomrule
\end{longtable}

\mcmd{MakiFitTikZ} measures its content in a savebox and only rescales it if the
content is wider than the requested width:

\begin{lstlisting}[language=TeX]
\sbox{\MakiFitTikZBox}{#2}
\ifdim\wd\MakiFitTikZBox>#1\relax
  \resizebox{#1}{!}{\usebox{\MakiFitTikZBox}}
\else
  \usebox{\MakiFitTikZBox}
\fi
\end{lstlisting}

This preserves the original size of small figures while preventing wide template
figures from overflowing the page.

\section{Box Environments and Counters}

Most content blocks are implemented with \mpkg{tcolorbox}. The style hierarchy is:

\begin{enumerate}
\item \mkey{commonbox}: shared appearance for theorem, definition, example, and exercise boxes.
\item \mkey{maki workflow card}/\mkey{maki workflow hero}: guide-card layouts.
\item \mkey{maki workflow block}: reusable block style for formulas, methods, research objects, and solutions.
\end{enumerate}

\begin{interfaceblock}{Main environments}
\item[\menv{theorem}, \menv{lemma}, \menv{proposition}, \menv{corollary}] Logical theorem-like blocks.
\item[\menv{definition}, \menv{axiom}] Definition-like blocks.
\item[\menv{example}, \menv{exercise}, \menv{realexam}] Example, exercise, and exam-problem blocks.
\item[\menv{notation}, \menv{question}, \menv{conjecture}, \menv{claim}] Research objects.
\item[\menv{keyformula}, \menv{pitfall}, \menv{methodnote}, \menv{examfocus}] Teaching workflow blocks.
\item[\menv{remark}, \menv{proofsketch}, \menv{idea}, \menv{obstacle}, \menv{nextstep}] Research-process blocks.
\end{interfaceblock}

\begin{implementation}{Shared lecture-object numbering}
Theorem, definition, example, and exercise boxes currently step the shared
\mkey{theorem} counter through this \mpkg{tcolorbox} option:
\begin{lstlisting}[language=TeX]
before title={\refstepcounter{theorem}}
\end{lstlisting}
This creates a single lecture object sequence. The source also declares
\mkey{definition} and \mkey{example} counters, but current box titles do not use
those counters for output.
\end{implementation}

\section{Solution Modes}

Solution behavior is parsed by \mfile{maki-notes.cls} and implemented by
\mfile{maki-notes.sty}.

\begin{interfaceblock}{Solution interface}
\item[\mopt{answers=false}] Default mode. Solution bodies are not printed.
\item[\mopt{answers=true,solutions=inline}] Solution boxes are printed in place.
\item[\mopt{answers=true,solutions=appendix}] Solution bodies are collected and printed by \mcmd{PrintSolutions}.
\item[\mopt{solutions=hidden}] Solutions are hidden even when \mopt{answers=true}.
\end{interfaceblock}

\begin{implementation}{Appendix-style solution collection}
\menv{solution} uses the \verb|+b| argument type of \mcmd{NewDocumentEnvironment}
to capture its full body. In appendix mode, the body is not typeset immediately;
it is stored in dynamic commands such as
\verb|\maki@solution@<n>@title| and \verb|\maki@solution@<n>@body|. Later,
\mcmd{PrintSolutions} iterates through the registered entries.
\end{implementation}

\begin{implementation}{Beamer compatibility}
Beamer already defines a theorem-like \menv{solution} environment. To avoid a
name collision, Maki's solution collection layer is only enabled outside Beamer
mode.
\end{implementation}

\section{Reference Helpers}

The package exposes lightweight Chinese reference wrappers:

\begin{lstlisting}[language=TeX]
\chapref{chap:intro}
\secref{sec:setup}
\thmref{thm:main}
\figref{fig:diagram}
\eqnref{eq:energy}
\end{lstlisting}

These helpers are stable wrappers around \mcmd{ref} and \mcmd{eqref}; they do not
load \mpkg{cleveref}. The benefit is low dependency risk and predictable output.
The tradeoff is that the user chooses the correct reference command.

\section{Beamer Theme Implementation}

\mfile{beamerthememaki.sty} implements three layouts:

\begin{longtable}{@{}p{0.22\linewidth}p{0.7\linewidth}@{}}
\toprule
Layout & Behavior \\
\midrule
\mopt{durham} & Default layout. Title page uses a theme-colored corner shape; headline shows horizontal section navigation; footline shows title and frame number with a progress bar.\\
\mopt{minimal} & Reduced decoration, small header/footline, better for dense technical slides.\\
\mopt{splitbar} & Stronger structural navigation, with a split color band and a distinctive outline page.\\
\bottomrule
\end{longtable}

Layout-specific templates are dispatched with \mcmd{csname}:

\begin{lstlisting}[language=TeX]
\setbeamertemplate{title page}{\csname MakiBeamerTitlePage@\MakiBeamerLayout\endcsname}
\setbeamertemplate{headline}{\csname MakiBeamerHeadline@\MakiBeamerLayout\endcsname}
\setbeamertemplate{footline}{\csname MakiBeamerFootline@\MakiBeamerLayout\endcsname}
\end{lstlisting}

\begin{implementation}{Theme reuse}
The Beamer theme consumes semantic colors such as \mkey{LogicColor},
\mkey{ExColor}, \mkey{DefColor}, and \mkey{MainText} from the shared style layer.
If the theme is loaded without that layer, it defines minimal fallback colors.
\end{implementation}

\section{Build, Test, and Release Automation}

The \mfile{Makefile} provides the main local entry points:

\begin{longtable}{@{}p{0.22\linewidth}p{0.7\linewidth}@{}}
\toprule
Target & Meaning \\
\midrule
\mkey{make main} & Build \mfile{document2.tex}.\\
\mkey{make example} & Build \mfile{example.tex}.\\
\mkey{make beamer-demo} & Build \mfile{beamer-demo.tex}.\\
\mkey{make tikz-templates} & Build \mfile{tikz-template-pages.tex}.\\
\mkey{make smoke} & Build the minimal packaging smoke document.\\
\mkey{make test} & Build the smoke document and all test documents.\\
\mkey{make all} & Build the main document, examples, Beamer demo, TikZ templates, and tests.\\
\bottomrule
\end{longtable}

\noindent\mfile{.github/workflows/test.yml}
builds smoke and test documents for push and pull request events.

\noindent\mfile{.github/workflows/release.yml}
builds the full example set, prepares release assets, and publishes tagged
releases.

\section{Extension Guidelines}

\subsection{Adding a Theme}

\begin{enumerate}
\item Add \mcmd{makiapplytheme<name>} in \mfile{maki-notes.sty}.
\item Define every semantic color slot, not only the colors that differ visibly.
\item Register the new value in README, template guide, and tests.
\item Use \mfile{tests/test-tikz-styles.tex} to verify that TikZ styles follow the theme.
\end{enumerate}

\subsection{Adding a Lecture Environment}

\begin{enumerate}
\item Decide whether the environment belongs to the shared \mkey{theorem} sequence.
\item Add or reuse a \mpkg{tcolorbox} style.
\item Expose the environment with \mcmd{NewDocumentEnvironment} or \mcmd{newenvironment}.
\item Add or extend a test document.
\end{enumerate}

\subsection{Adding a TikZ Style Family}

\begin{enumerate}
\item Use semantic names such as \mkey{flow process} or \mkey{graph edge strong}.
\item Reference semantic color slots instead of fixed palette values.
\item Add a copyable example page to \mfile{tikz-template-pages.tex}.
\item Add minimal coverage to \mfile{tests/test-tikz-styles.tex}.
\end{enumerate}

\section{Maintenance Notes}

\begin{itemize}
\item Keep the \mfile{.cls} files thin; large visual logic belongs in style files.
\item Parse new public options in the relevant class, then pass resolved values to the style layer.
\item Check \mcmd{ifMakiBeamerMode} whenever a feature might touch Beamer.
\item Features involving floats, captions, wrap figures, marginalia, or page headers normally belong to notes mode only.
\item New counters and indexes should define label/reference behavior and test coverage.
\item Run \verb|make -B test| after interface changes; run \verb|make -B all| before release.
\end{itemize}

\end{document}
